% Generated draft: requires mathematical and source review.
\documentclass[11pt]{article}
\usepackage[margin=25mm]{geometry}
\usepackage{fontspec}
\setmainfont{Latin Modern Roman}
\usepackage{amsmath,amssymb,mathtools,unicode-math}
\setmathfont{Latin Modern Math}
\usepackage{enumitem,needspace,xurl,hyperref}
\setlength{\parindent}{0pt}
\setlength{\parskip}{5pt}
\setlength{\emergencystretch}{3em}
\newcommand{\N}{\mathbb{N}}
\newcommand{\Z}{\mathbb{Z}}
\newcommand{\Q}{\mathbb{Q}}
\newcommand{\R}{\mathbb{R}}
\newcommand{\C}{\mathbb{C}}
\newcommand{\F}{\mathbb{F}}
\newcommand{\A}{\mathbb{A}}
\newcommand{\PP}{\mathbb{P}}
\newcommand{\E}{\mathbb{E}}
\newcommand{\eps}{\varepsilon}
\newcommand{\dd}{\,\mathrm{d}}
\begin{document}
\section*{UnsolvedMath ID: 30006029}
\textbf{Source problem code:} catalogue entry
\subsection{Definitions and mathematical statement}
For a smooth, connected, compact complex curve \(C\) of genus \(g\), a meromorphic \(k\)-differential, in the usual convention, is a nonzero meromorphic section \(\eta\) of \(K_C^{\otimes k}\), where \(k\in\Z_{>0}\) and \(K_C\) is the canonical line bundle. If its complete zero--pole divisor is supported at distinct marked points \(p_1,\ldots,p_n\), write
\[
\operatorname{div}(\eta)=\sum_{i=1}^{n}m_i p_i,
\qquad m_i\in\Z,
\qquad \sum_{i=1}^{n}m_i=k(2g-2).
\]
Positive orders describe zeros and negative orders describe poles. This explains the ordinary meaning of a zero--pole profile; the supplied record does not specify whether its parameter \(A\) records these orders directly or uses a shifted convention. No identification of \(A\) with \((m_1,\ldots,m_n)\) is therefore asserted.

\textbf{Problem.} For positive \(k\), arbitrary genus \(g\), and each admissible profile \(A\) in the convention of the cited source, find an enumerative interpretation of the \(k\)-leaky double Hurwitz number
\[
H_g^k(A).
\]
More broadly, formulate natural enumerative problems for meromorphic \(k\)-differentials of genus \(g\) and zero--pole profile \(A\), and determine whether their counts equal, or have another explicitly specified relationship with, \(H_g^k(A)\). An interpretation must specify the objects counted, the conditions making the count meaningful, and any multiplicities or automorphism weights. The supplied record does not define \(H_g^k(A)\), its normalization, its admissible profiles, or the relevant compactification and boundary conditions. Thus this is a research question referring to the source's invariant, rather than a self-contained numerical identity.

\textbf{Supplied motivation and unresolved qualifications.} The excerpt mentions a vector \(e=(e_1,\ldots,e_n)\) of nonnegative integers and a variant \(b_{\mathrm{reg}}(A)\) of a branch class having codimension \(2g-3+n-|e|\), with the usual interpretation \(|e|=\sum_i e_i\). It describes unspecified numbers depending on this data as interpolating between double Hurwitz numbers at \(e=0\) and intersection numbers involving the usual double ramification cycle when \(e\) is maximal. It also reports piecewise polynomiality, wall-crossing formulas, and expressions as sums over tropical covers with vertex leaks reflecting discrepancies between zero and pole multiplicities of \(k\)-differentials. The defining formula, the meaning of maximal \(e\), and the precise relationship of these numbers to \(H_g^k(A)\) are absent; these descriptions are retained as supplied background, not additional established hypotheses.

The excerpt further discusses a genus-zero locus with corrupted double-ramification notation, carrying residue conditions at \(p_2,\ldots,p_{n+1}\). It describes the residues of \(\eta\) as being ``linearly dependent to'' a fixed vector
\[
\widetilde r=(r_2,\ldots,r_{n+1})\in\C^n,
\qquad \sum_{i=2}^{n+1}r_i=0.
\]
For generic \(\widetilde r\), it reports the cardinality
\[
(d-1)(d-2)\cdots\bigl(d-(n-2)\bigr)
\tag{10}
\]
for that locus, referring to unidentified references [6, 2, 5]. The meanings and ranges of \(d,n\), the profile in this example, the equivalence relation on differentials, and the intended residue condition are not supplied. The accompanying formula (11) for the genus-zero \(1\)-leaky number \(H_0^1(A)\) is too corrupted to reconstruct reliably. The asserted structural similarity of these two formulas is motivation for the question; no equality or relationship is asserted here.
\subsection{Short English statement}
Seek a geometric counting interpretation of the positive-\(k\) leaky double Hurwitz numbers \(H_g^k(A)\), and natural counting problems for \(k\)-differentials in arbitrary genus whose answers relate to these numbers. The supplied source leaves the invariant's definition, profile convention, counting weights, and boundary conditions unspecified, and its motivating residue example is partly corrupted. This statement-only draft requires mathematical review.
\subsection{Sources}
\begin{itemize}
\item Supplied local UnsolvedMath registry record, permanent ID \(30006029\), source problem code \texttt{catalogue entry}, titled \emph{Enumerative Meaning of k-Leaky Double Hurwitz Numbers}; collection: Oberwolfach Reports Open Problems. No individual proposer is supplied. The registry records proposed year 2024.
\item Supplied source citation: \emph{Riemann Surfaces: Random, Flat, and Hyperbolic Geometry} (2025), Oberwolfach Reports, pp.\ 2369--2424, DOI \texttt{10.4171/owr/2024/41}. Source URL and literature-source link labeled ``OWR report'': \url{https://doi.org/10.4171/owr/2024/41}. The supplied citation year 2025 and proposed year 2024 are preserved without resolving their relationship.
\item Supplied literature-source attribution: Cavalieri--Markwig--Schmitt, \url{https://doi.org/10.1017/fms.2025.26}. No title or theorem-level identification is supplied.
\item Supplied literature-source label: ``one-part leaky covers'', \url{https://arxiv.org/abs/2509.04335}. No authors or theorem-level identification are supplied.
\item Literature assessment checked at \texttt{2026-08-22}: the supplied curation labels the problem open and states that intersection-theoretic and tropical formulas provide structural information while the positive-\(k\) enumerative meaning remains unclear. This is supplied dated curation, not an independently verified literature assessment. The record explicitly requests independent verification of time-sensitive status and cited claims.
\item The registry labels the statement \texttt{corrected\_verified} and says that only formula delimiters, notation, grammar, and typography were repaired while preserving mathematical content and scope. This supplied label is not independent mathematical verification. The excerpt's numerical references [6, 2, 5] have no supplied bibliography mapping.
\end{itemize}
\end{document}
